%!TEX root = ../nauchka.tex

\section{Эксперименты}

\begin{frame}
\frametitle{Аппаратная конфигурация и план}

\textbf{Аппаратура:} 2$\times$ NVIDIA A40 (48~ГБ), NVLink, NCCL, PyTorch~2.x

\vspace{0.3cm}

\textbf{5 групп экспериментов:}
\begin{enumerate}
    \item TP: потребление GPU-памяти (OOM-тест)
    \item TP: корректность и soundness bounds
    \item TP: автоматическое шардирование моделей VNN-COMP
    \item FSDP: корректность bounds
    \item FSDP: потребление GPU-памяти
\end{enumerate}

\end{frame}

\begin{frame}
\frametitle{Эксперимент 1: TP --- потребление памяти}

\textbf{Модель:} MLP, $D{=}4096$, $H{=}262144$, batch$=$2048, CROWN, $\varepsilon{=}0.01$

\vspace{0.3cm}

\begin{table}[h]
\centering
\scriptsize
\begin{tabular}{lccc}
\hline
\textbf{Режим} & \textbf{GPU} & \textbf{max\_allocated, МБ} & \textbf{max\_reserved, МБ} \\
\hline
Single GPU & 1$\times$A40 & 26\,826 & 30\,874 \\
TP = 2     & GPU 0        & 13\,513 & 15\,512 \\
TP = 2     & GPU 1        & 13\,513 & 15\,512 \\
\hline
\end{tabular}
\end{table}

\vspace{0.3cm}

\begin{exampleblock}{Результат}
Коэффициент снижения: $26826 / 13513 \approx \mathbf{1.985\times}$ (теор. предел $2.0\times$)
\end{exampleblock}

\end{frame}

\begin{frame}
\frametitle{Эксперимент 2: TP --- корректность bounds}

\textbf{Метод:} CROWN, $\varepsilon{=}0.02$; сравнение TP$=$2 с single-GPU эталоном

\vspace{0.2cm}

\begin{table}[h]
\centering
\scriptsize
\begin{tabular}{lccc}
\hline
\textbf{Модель} & \textbf{Пар Col-Row} & $|lb_{\text{diff}}|$ & $|ub_{\text{diff}}|$ \\
\hline
PyTorch 256$\times$2 & 1 & $3 \cdot 10^{-8}$ & $3 \cdot 10^{-8}$ \\
PyTorch 256$\times$4 & 2 & $2.5 \cdot 10^{0}$  & $2.5 \cdot 10^{0}$ \\
ONNX 256$\times$2    & 1 & $6 \cdot 10^{-8}$ & $6 \cdot 10^{-8}$ \\
ONNX 256$\times$4    & 2 & $5.6 \cdot 10^{0}$  & $4.0 \cdot 10^{0}$ \\
ONNX 256$\times$6    & 3 & $7.9 \cdot 10^{2}$  & $7.5 \cdot 10^{2}$ \\
\hline
\end{tabular}
\end{table}

\vspace{0.2cm}

\begin{alertblock}{Вывод}
Soundness подтверждена для всех конфигураций ($lb_{\text{tp}} \leq lb_{\text{ref}}$, $ub_{\text{tp}} \geq ub_{\text{ref}}$).
\begin{itemize}
    \item 1 пара Col-Row: совпадение с точностью $\sim 10^{-8}$ (машинный $\varepsilon$)
    \item 2+ пар: bounds шире эталона, ошибка растёт на порядки (IBP fallback)
\end{itemize}
\end{alertblock}

\end{frame}

\begin{frame}
\frametitle{Эксперимент 3: автошардирование VNN-COMP моделей}

\begin{block}{Задача}
Проверить \texttt{tp\_shard\_bounded\_module} на ONNX-моделях MNIST-FC из VNN-COMP 2022
\end{block}

\vspace{0.3cm}

\textbf{Процесс:}
\begin{enumerate}
    \item Загрузка ONNX $\to$ PyTorch (\texttt{onnx2pytorch})
    \item Обёртка в \texttt{BoundedModule}
    \item Автоматическое шардирование без ручного вмешательства
\end{enumerate}

\vspace{0.3cm}

\begin{exampleblock}{Результат}
Все 3 модели (256$\times$2, 256$\times$4, 256$\times$6) успешно шардированы и прошли тест на soundness
\end{exampleblock}

\end{frame}

\begin{frame}
\frametitle{Эксперимент 4: FSDP --- корректность bounds}

\textbf{Метод:} IBP + CROWN, FSDP$=$2 vs single-GPU; 8 конфигураций

\vspace{0.2cm}

\begin{table}[h]
\centering
\scriptsize
\begin{tabular}{llcc}
\hline
\textbf{Модель} & \textbf{Метод} & $|\Delta lb|$ & $|\Delta ub|$ \\
\hline
MLP 256$\times$2 & IBP   & $0.00$ & $0.00$ \\
MLP 256$\times$2 & CROWN & $0.00$ & $0.00$ \\
MLP 256$\times$4 & IBP   & $0.00$ & $0.00$ \\
MLP 256$\times$4 & CROWN & $0.00$ & $0.00$ \\
MLP 256$\times$6 & CROWN & $0.00$ & $0.00$ \\
ONNX 256$\times$2 & CROWN & $0.00$ & $0.00$ \\
ONNX 256$\times$4 & CROWN & $0.00$ & $0.00$ \\
\hline
\end{tabular}
\end{table}

\vspace{0.2cm}

\begin{exampleblock}{Результат}
\textbf{Побитово идентичные} bounds для всех моделей и методов ($\Delta = 0.00$, не $\sim\!10^{-8}$). Кросс-ранговое расхождение: $0.00$.
\end{exampleblock}

\end{frame}

\begin{frame}
\frametitle{Эксперимент 5: FSDP --- потребление памяти}

\begin{columns}[T]
\begin{column}{0.48\textwidth}
\textbf{Baseline-память (веса):}

\scriptsize
\begin{tabular}{lrrr}
\hline
\textbf{Модель} & \textbf{Single} & \textbf{FSDP} & \textbf{$\downarrow$\%} \\
\hline
$h{=}256, d{=}4$  & 26 & 11 & 57 \\
$h{=}1024, d{=}4$ & 187 & 40 & 79 \\
$h{=}4096, d{=}4$ & 2272 & 417 & 82 \\
$h{=}8192, d{=}4$ & 8760 & 1594 & 82 \\
$h{=}4096, d{=}8$ & 9721 & 930 & \textbf{90} \\
\hline
\end{tabular}
\normalsize
\end{column}

\begin{column}{0.48\textwidth}
\textbf{Пиковая память:}

\scriptsize
\begin{tabular}{lrrr}
\hline
\textbf{Модель} & \textbf{Single} & \textbf{FSDP} & \textbf{$\downarrow$\%} \\
\hline
$h{=}4096, d{=}4$  & 3800 & 2527 & 34 \\
$h{=}8192, d{=}4$  & 14759 & 9782 & 34 \\
$h{=}4096, d{=}8$  & 16726 & 10232 & \textbf{39} \\
\hline
\end{tabular}
\normalsize
\end{column}
\end{columns}

\vspace{0.3cm}

Размеры в МБ. Теория: $M_W^{\text{FSDP}}/M_W^{\text{single}} = 1/P + 1/d$

\end{frame}

\begin{frame}
\frametitle{Сравнение TP и FSDP}

\begin{table}[h]
\centering
\small
\begin{tabular}{|l|c|c|}
\hline
& \textbf{TP} & \textbf{FSDP} \\
\hline
Пиковая память & $\sim\!2\times$ снижение & 34--39\% снижение \\
\hline
Baseline-память & $\sim\!2\times$ снижение & 80--90\% снижение \\
\hline
Точность bounds & Падает с глубиной & Побитово идентична \\
\hline
Совместимость с $\beta$-CROWN & Сложно & Да \\
\hline
Объём реализации & $\sim$500 строк & $\sim$100 строк \\
\hline
\end{tabular}
\end{table}

\vspace{0.3cm}

\begin{block}{Вывод}
TP --- максимальная экономия памяти для неглубоких моделей; \\
FSDP --- идеальная точность и совместимость с полной верификацией
\end{block}

\end{frame}